% ch-a4-14.tex —— 第 4 篇（Emerton–Gee–Hellmann）续片，底本 PDF 336–347 = 印刷页 323–334

\ysrc{323}
虽然 $U$ 与 $U^s$ 作为 $\mathrm{GL}_2(\F_{p^2})$ 的表示不同构，但它们到 $\mathrm{GL}_2(\Q_{p^2})$ 的紧诱导同构（经由同构 (7.7.8)），因此我们同样有
\[
\mathfrak A\bigl(c\text{-}\mathrm{Ind}_{KZ}^GU^s\bigr)=\OO_{C_0\cup C_1\cup C_2\cup C_3}(-1,-1,0,0,0).
\]
注意：若以 $(-1,-1,0,0,0)$ 对图 (7.7.28) 作扭，便得到图
\begin{equation}\tag{7.7.35}
\begin{tikzcd}[column sep=large,row sep=large]
& \OO_{C_1}(-1,0)\arrow[dl,no head]\arrow[dr,no head] & \\
\OO_{C_0}(-1,0)\arrow[dr,no head] & & \OO_{C_2}(-1,0)\arrow[dl,no head]\\
& \OO_{C_3}(-1,0) &
\end{tikzcd}
\end{equation}
这验证了 $\mathfrak A$ 作用于图 (7.7.34) 时的函子性。

复合映射
\[
c\text{-}\mathrm{Ind}_{KZ}^G\sigma_1\ra c\text{-}\mathrm{Ind}_{KZ}^GU\isom c\text{-}\mathrm{Ind}_{KZ}^GU^s\ra c\text{-}\mathrm{Ind}_{KZ}^G\sigma_1
\]
已知与 Hecke 算子 $T$ 重合。考察例 7.7.26 中的诸图可知，该映射由乘以 $v_1y_1$ 给出，而这正是 $\mathfrak A(T)$。与此计算相关，我们指出：(7.7.21) 与 $\mathfrak A$ 的全忠实性、以及 $c\text{-}\mathrm{Ind}_{KZ}^GU$ 的自同态环的已知结构相容。

众所周知，通过约化整主级数中的格，我们还能得到 $\mathrm{GL}_2(\F_q)$ 的表示 $V$ 与 $W$，它们可分别用扩展图（extension diagram）
\begin{equation}\tag{7.7.36}
\begin{tikzcd}[column sep=large,row sep=large]
& \sigma_0\arrow[dl,no head]\arrow[dr,no head] & \\
\sigma_3\arrow[dr,no head] & & \sigma_1\arrow[dl,no head]\\
& \sigma_2 &
\end{tikzcd}
\end{equation}
与
\begin{equation}\tag{7.7.37}
\begin{tikzcd}[column sep=large,row sep=large]
& \sigma_2\arrow[dl,no head]\arrow[dr,no head] & \\
\sigma_1\arrow[dr,no head] & & \sigma_3\arrow[dl,no head]\\
& \sigma_0 &
\end{tikzcd}
\end{equation}
描述。我们定义
\[
\mathfrak A\bigl(c\text{-}\mathrm{Ind}_{KZ}^GV\bigr):=\OO_{C_0\cup C_1\cup C_2\cup C_3}(-1,-1,-1,0,0)
\]
以及
\[
\mathfrak A\bigl(c\text{-}\mathrm{Ind}_{KZ}^GW\bigr):=\OO_{C_0\cup C_1\cup C_2\cup C_3}(-1,-1,1,0,0).
\]

\ysrc{324}
然后分别以 $(-1,-1,-1,0,0)$ 与 $(-1,-1,1,0,0)$ 对图 (7.7.30) 与 (7.7.31) 作扭，得到图
\begin{equation}\tag{7.7.38}
\begin{tikzcd}[column sep=large,row sep=large]
& \OO_{C_0}(-1,0)\arrow[dl,no head]\arrow[dr,no head] & \\
\OO_{C_3}(-1,0)\arrow[dr,no head] & & \OO_{C_1}(-1,0)\arrow[dl,no head]\\
& \OO_{C_2}(-1,0) &
\end{tikzcd}
\end{equation}
与
\begin{equation}\tag{7.7.39}
\begin{tikzcd}[column sep=large,row sep=large]
& \OO_{C_2}(-1,0)\arrow[dl,no head]\arrow[dr,no head] & \\
\OO_{C_1}(-1,0)\arrow[dr,no head] & & \OO_{C_3}(-1,0)\arrow[dl,no head]\\
& \OO_{C_0}(-1,0) &
\end{tikzcd}
\end{equation}
这验证了 $\mathfrak A$ 作用于 (7.7.36) 与 (7.7.37) 时的函子性。

\subsubsection*{7.7.40.\quad 四分支之并——情形二（尖点）}
此情形的计算与主级数情形的计算类似。不过，对驯顺尖点型（tame cuspidal type）而言，并不存在优先的格。其约化（一般地）含四个权 $\sigma_0$、$\sigma_1$、$\sigma_2$ 与 $\sigma_3$，每个重数一，并且可以安排它们出现在扩展图
\[
\begin{tikzcd}[column sep=large,row sep=large]
& \sigma_0\arrow[dl,no head]\arrow[dr,no head] & \\
\sigma_3\arrow[dr,no head] & & \sigma_1\arrow[dl,no head]\\
& \sigma_2 &
\end{tikzcd}
\]
中，更一般地出现在该图的任一循环置换中。

在对权作适当标号之后，并 $C_{\sigma_0}\cup C_{\sigma_1}\cup C_{\sigma_2}\cup C_{\sigma_3}$ 容许开覆盖
\begin{equation}\tag{7.7.41}
\begin{aligned}
\bigl\{&D(v_0x_0),\ D(u_1y_1),\ D(v_2x_2),\ D(u_3y_3),\ D(v_0y_0)\cup D(u_1x_1),\\
&D(v_1y_1)\cup D(u_2x_2),\ D(v_2y_2)\cup D(u_3x_3),\ D(v_3y_3)\cup D(u_0x_0),\\
&D(u_0y_0)\cup D(v_1x_1)\cup D(u_2y_2)\cup D(v_3x_3)\bigr\}.
\end{aligned}
\end{equation}
我们有
\[
\begin{aligned}
D(u_iy_i)&=\mathbf A^1_{x_i/u_i}\times\mathbf A^1_{v_iy_i},\\
D(v_ix_i)&=\mathbf A^1_{u_i/x_i}\times\mathbf A^1_{v_iy_i},\\
D(u_ix_i)&=(\mathbf G_m)_{x_i/u_i}\times\bigl[\mathbf A^2_{v_i,y_i}/(\mathbf G_m)_{b_i}\bigr],\\
D(v_iy_i)&=\bigl[\mathbf A^2_{u_i,x_i}/(\mathbf G_m)_{a_i}\bigr]\times(\mathbf G_m)_{v_iy_i},
\end{aligned}
\]

\ysrc{325}
且对 $i=0,2$，有
\[
D(v_iy_i)\cup D(u_{i+1}x_{i+1})=\Bigl[\bigl\{(u_i,x_i,v_iy_i,y_{i+1})\bigm| v_iy_i\neq0,\ x_iy_{i+1}=0\bigr\}\big/(\mathbf G_m)_{a_i}\Bigr],
\]
其中等同关系为 $u_i=v_{i+1}$、$v_iy_i=u_{i+1}/x_{i+1}$ 与 $a_i=b_{i+1}$；而对 $i=1,3$，有
\[
D(v_iy_i)\cup D(u_{i+1}x_{i+1})=\Bigl[\bigl\{(u_i,x_i,v_iy_i,v_{i+1})\bigm| v_iy_i\neq0,\ u_iv_{i+1}=0\bigr\}\big/(\mathbf G_m)_{a_i}\Bigr],
\]
其中等同关系为 $x_i=y_{i+1}$、$v_iy_i=x_{i+1}/u_{i+1}$ 与 $a_i=b_{i+1}^{-1}$；最后
\[
\begin{aligned}
&D(u_0y_0)\cup D(v_1x_1)\cup D(u_2y_2)\cup D(v_3x_3)\\
&\qquad=\bigl\{(v_0y_0,v_1y_1,v_2y_2,v_3y_3)\bigm|(v_0y_0)(v_2y_2)=(v_1y_1)(v_3y_2)=0\bigr\},
\end{aligned}
\]
其中等同关系为 $x_0/u_0=v_3y_3$、$u_1/x_1=v_0y_0$、$x_2/u_2=v_1y_1$ 与 $u_3/x_3=v_2y_2$\tn{底本此处作 $(v_1y_1)(v_3y_2)=0$；按本段末尾的等同关系 $x_0/u_0=v_3y_3$ （以及 $u_3/x_3=v_2y_2$）所体现的模式，第二项应为 $v_3y_3$，疑为底本笔误。译文照录原文。}。

利用开覆盖 (7.7.41)，我们计算得
\begin{equation}\tag{7.7.42}
\Gamma\bigl(\OO_{C_0\cup C_1\cup C_2\cup C_3}\bigr)=k.
\end{equation}
我们还计算得
\[
\mathrm{Pic}\bigl(C_0\cup C_1\cup C_2\cup C_3\bigr)=\Z\times\Z\times\Z\times\Z,
\]
其中各因子依次由 $\mathrm{Pic}\bigl(D(v_2y_2)\cup D(u_3x_3)\bigr)$、$\mathrm{Pic}\bigl(D(v_1y_1)\cup D(u_2x_2)\bigr)$、$\mathrm{Pic}\bigl(D(v_0y_0)\cup D(u_1x_1)\bigr)$ 与 $\mathrm{Pic}\bigl(D(v_3y_3)\cup D(u_0x_0)\bigr)$ 给出。我们通过描述嵌入
\begin{equation}\tag{7.7.43}
\mathrm{Pic}\bigl(C_0\cup C_1\cup C_2\cup C_3\bigr)\into\mathrm{Pic}(C_0)\times\mathrm{Pic}(C_1)\times\mathrm{Pic}(C_2)\times\mathrm{Pic}(C_3)
\end{equation}
来固定坐标的选取；该嵌入由
\begin{equation}\tag{7.7.44}
(k,l,m,n)\mapsto\bigl((m,n),(-l,m),(k,l),(-n,k)\bigr)
\end{equation}
给出。

\begin{exam}{7.7.45}\label{exam:a4-7-7-45}
我们有下列 $\OO_{C_0\cup C_1\cup C_2\cup C_3}$ 的理想子层（ideal subsheaf）图：
{\small
\[
\begin{tikzcd}[column sep=tiny,row sep=large]
 & \OO_{C_0\cup C_1\cup C_2\cup C_3} & & \\
 & \mathcal I_{C_0}\arrow[u,"\OO_{C_0}"]\arrow[r,equal,no head] & \mathcal I_{C_0\cup C_3}+\mathcal I_{C_0\cup C_1} & \\
 \mathcal I_{C_0\cup C_3}\arrow[ur,hook,"{\OO_{C_3}(-1,0)}"] & & \mathcal I_{C_0\cup C_3}\arrow[ul,hook,"{\OO_{C_1}(0,1)}"'] & \\
 & \mathcal I_{C_0\cup C_1\cup C_3}\arrow[ul,hook]\arrow[rr,equal,no head] & & \mathcal I_{C_0\cup C_3}\cap\mathcal I_{C_0\cup C_1}\arrow[ul,hook] & & \OO_{C_2}(-1,-1)\arrow[ll,"\sim"]
\end{tikzcd}
\]}
\tn{底本此图右中与左中的两个节点都印作 $\mathcal I_{C_0\cup C_3}$；按本图其余三幅（以及下方求交的写法 $\mathcal I_{C_0\cup C_3}\cap\mathcal I_{C_0\cup C_1}$）的模式，右中节点应为 $\mathcal I_{C_0\cup C_1}$。译文照录原文。}
\ysrc{326}
{\small
\[
\begin{tikzcd}[column sep=tiny,row sep=large]
 & \OO_{C_0\cup C_1\cup C_2\cup C_3} & & \\
 & \mathcal I_{C_1}\arrow[u,"\OO_{C_1}"]\arrow[r,equal,no head] & \mathcal I_{C_0\cup C_1}+\mathcal I_{C_1\cup C_2} & \\
 \mathcal I_{C_0\cup C_1}\arrow[ur,hook,"{\OO_{C_0}(-1,0)}"] & & \mathcal I_{C_1\cup C_2}\arrow[ul,hook,"{\OO_{C_2}(0,-1)}"'] & \\
 & \mathcal I_{C_0\cup C_1\cup C_2}\arrow[ul,hook]\arrow[rr,equal,no head] & & \mathcal I_{C_0\cup C_1}\cap\mathcal I_{C_1\cup C_2}\arrow[ul,hook] & & \OO_{C_3}(-1,1)\arrow[ll,"\sim"]
\end{tikzcd}
\]}
{\small
\[
\begin{tikzcd}[column sep=tiny,row sep=large]
 & \OO_{C_0\cup C_1\cup C_2\cup C_3} & & \\
 & \mathcal I_{C_2}\arrow[u,"\OO_{C_2}"]\arrow[r,equal,no head] & \mathcal I_{C_1\cup C_2}+\mathcal I_{C_2\cup C_3} & \\
 \mathcal I_{C_1\cup C_2}\arrow[ur,hook,"{\OO_{C_1}(-1,0)}"] & & \mathcal I_{C_2\cup C_3}\arrow[ul,hook,"{\OO_{C_3}(0,1)}"'] & \\
 & \mathcal I_{C_1\cup C_2\cup C_3}\arrow[ul,hook]\arrow[rr,equal,no head] & & \mathcal I_{C_1\cup C_2}\cap\mathcal I_{C_2\cup C_3}\arrow[ul,hook] & & \OO_{C_0}(-1,-1)\arrow[ll,"\sim"]
\end{tikzcd}
\]}
{\small
\[
\begin{tikzcd}[column sep=tiny,row sep=large]
 & \OO_{C_0\cup C_1\cup C_2\cup C_3} & & \\
 & \mathcal I_{C_3}\arrow[u,"\OO_{C_3}"]\arrow[r,equal,no head] & \mathcal I_{C_2\cup C_3}+\mathcal I_{C_0\cup C_3} & \\
 \mathcal I_{C_2\cup C_3}\arrow[ur,hook,"{\OO_{C_2}(-1,0)}"] & & \mathcal I_{C_0\cup C_3}\arrow[ul,hook,"{\OO_{C_0}(0,-1)}"'] & \\
 & \mathcal I_{C_0\cup C_2\cup C_3}\arrow[ul,hook]\arrow[rr,equal,no head] & & \mathcal I_{C_2\cup C_3}\cap\mathcal I_{C_0\cup C_3}\arrow[ul,hook] & & \OO_{C_1}(-1,1)\arrow[ll,"\sim"]
\end{tikzcd}
\]}
\end{exam}

下列表示（用其扩展图描述）作为尖点型中格的约化出现：
\begin{equation}\tag{7.7.46}
\begin{tikzcd}[column sep=large,row sep=large]
& \sigma_0\arrow[dl,no head]\arrow[dr,no head] & \\
\sigma_3\arrow[dr,no head] & & \sigma_1\arrow[dl,no head]\\
& \sigma_2 &
\end{tikzcd}
\end{equation}

\ysrc{327}
\begin{equation}\tag{7.7.47}
\begin{tikzcd}[column sep=large,row sep=large]
& \sigma_1\arrow[dl,no head]\arrow[dr,no head] & \\
\sigma_0\arrow[dr,no head] & & \sigma_2\arrow[dl,no head]\\
& \sigma_3 &
\end{tikzcd}
\end{equation}
\begin{equation}\tag{7.7.48}
\begin{tikzcd}[column sep=large,row sep=large]
& \sigma_2\arrow[dl,no head]\arrow[dr,no head] & \\
\sigma_1\arrow[dr,no head] & & \sigma_3\arrow[dl,no head]\\
& \sigma_0 &
\end{tikzcd}
\end{equation}
以及
\begin{equation}\tag{7.7.49}
\begin{tikzcd}[column sep=large,row sep=large]
& \sigma_3\arrow[dl,no head]\arrow[dr,no head] & \\
\sigma_2\arrow[dr,no head] & & \sigma_0\arrow[dl,no head]\\
& \sigma_1 &
\end{tikzcd}
\end{equation}
相应的 $\mathfrak A$ 取值为
\[
\begin{aligned}
&\OO_{C_0\cup C_1\cup C_2\cup C_3}(0,1,-1,0),\quad
 \OO_{C_0\cup C_1\cup C_2\cup C_3}(-1,1,0,0),\\
&\OO_{C_0\cup C_1\cup C_2\cup C_3}(-1,0,0,1)\quad\text{与}\quad
 \OO_{C_0\cup C_1\cup C_2\cup C_3}(0,0,-1,1)\text{。}
\end{aligned}
\]若以相应的扭对例 7.7.45 中的每个图作扭（并用公式 (7.7.44) 计算），我们就得到一个由诸 $\mathfrak A(\sigma_i)$ 按适当次序拼成的扩展图；这依次验证了 $\mathfrak A$ 在诸图 (7.7.46)、(7.7.47)、(7.7.48) 与 (7.7.49) 上的函子性。

\subsubsection*{7.7.50.\quad 一个非紧诱导的例子}
同构 $c\text{-}\mathrm{Ind}_{KZ}^G\ol{\mathrm{PS}}(\xi)\isom c\text{-}\mathrm{Ind}_{KZ}^G\ol{\mathrm{PS}}(\xi^s)$ 使我们能构造一个有趣的表示。（这里 $\ol{\mathrm{PS}}(\xi)$ 就是 §7.7.19 中记作 $U$ 的表示，$\ol{\mathrm{PS}}(\xi^s)$ 是记作 $U^s$ 的表示。）若以 $\sigma_\xi$ 记 $\ol{\mathrm{PS}}(\xi)$ 的基座（socle），以 $\sigma_{\xi^s}$ 记 $\ol{\mathrm{PS}}(\xi^s)$ 的基座，则有嵌入
\[
c\text{-}\mathrm{Ind}_{KZ}^G\sigma_\xi\into c\text{-}\mathrm{Ind}_{KZ}^G\ol{\mathrm{PS}}(\xi)
\quad\text{与}\quad
c\text{-}\mathrm{Ind}_{KZ}^G\sigma_{\xi^s}\into c\text{-}\mathrm{Ind}_{KZ}^G\ol{\mathrm{PS}}(\xi^s)\isom c\text{-}\mathrm{Ind}_{KZ}^G\ol{\mathrm{PS}}(\xi),
\]
从而诱导出
\[
c\text{-}\mathrm{Ind}_{KZ}^G\sigma_\xi\oplus c\text{-}\mathrm{Ind}_{KZ}^G\sigma_{\xi^s}\into c\text{-}\mathrm{Ind}_{KZ}^G\ol{\mathrm{PS}}(\xi).
\]
定义 $\Pi$ 为该嵌入的余核。

利用 §7.7.19 的计算，我们可以算出 $\mathfrak A(\Pi)$。事实上，注意在该讨论的记号下 $\sigma_\xi$ 与 $\sigma_{\xi^s}$ 可以取为 $\sigma_1$ 与 $\sigma_3$，我们发现 $\mathfrak A(\Pi)$ 是嵌入
\[
\OO_{C_1}(-1,0)\oplus\OO_{C_3}(-1,0)\into\OO_{C_0\cup C_1\cup C_2\cup C_3}(-1,-1,0,0,0)
\]
的余核。

\ysrc{328}
该嵌入的像可等同于理想子层 $\mathcal I_{C_0\cup C_2}$ 经 $(-1,-1,0,0,0)$ 的扭，并且可以算出 $\mathfrak A(\Pi)$ 是 $C_0\cup C_2$ 上的某个可逆层。（特别地，它不是 Cohen--Macaulay 的。）

\subsubsection*{7.7.51.\quad 一个例子——权循环}
在 2006 年的 AIM 会议上，与会者找到了一个论证，它后来被称为「权循环」（weight cycling）。其原始形式如下。（欲知更多细节，感兴趣的读者可参看当时的报告 \url{https://aimath.org/WWN/padicmodularity/emertonreport.pdf} 与 \url{https://www.ma.ic.ac.uk/~buzzard/maths/research/notes/notes_on_mod_p_local_langlands.pdf}。）

\begin{prop}{7.7.52}\label{prop:a4-7-7-52}
设 $\pi'$ 是 $\mathrm{GL}_2(\Q_{p^2})$ 的表示，满足下列性质：
\begin{enumerate}[(1)]
\item $\pi'$ 的 $K$-基座（$K$-socle）的各组成因子（constituent）都在某个充分一般的尖点型（cuspidal type）$\tau$ 的约化之中；
\item $\pi'$ 的 $K$-基座的每个组成因子都以重数一出现；
\item $\pi'$ 的 $K$-基座的每个组成因子都是超奇异的。
\end{enumerate}
则 $\pi'$ 有一个（必然不可约的）子表示 $\pi$，具有下列性质：
\begin{enumerate}[(a)]
\item $\pi$ 有中心特征（central character）；
\item $\pi'$ 的每个 $G$-子表示都含 $\pi$；
\item $\pi$ 的 $K$-基座与 $\pi'$ 的 $K$-基座相同；
\item $\tau$ 的约化的四个组成因子都实际出现在 $\pi'$ 与 $\pi$ 的公共 $K$-基座中。
\end{enumerate}
\end{prop}

\begin{proof}
$\pi'$ 的任一 $G$-子表示都至少含 $\pi'$ 的 $K$-基座的一个组成因子，从而含该组成因子生成的 $G$-子表示。因此，只要能证明：由其中任一组成因子生成的 $G$-子表示 $\pi$ 含所有这些组成因子——也就是说，只要能证明 $\pi$ 满足 (c)——我们也就证明了 $\pi$ 与初始所选取的生成组成因子无关，从而 $\pi$ 也满足条件 (b)。（而 $\pi$ 满足 (b) 这一事实蕴含它不可约。）

于是我们转而证明 $\pi$ 确实满足 (c)。在此过程中，我们还将证明它满足 (a) 与 (d)。

以 $\sigma_0$ 记 $\tau$ 的约化中出现的那个选定的组成因子，它出现在 $\pi'$ 的 $K$-基座中并生成 $\pi$。由于 $\sigma_0$ 在 $\pi$ 的 $K$-基座中以重数一出现，$G$ 的中心 $Z$ 通过某个特征作用在 $\sigma_0$ 上，从而也通过该特征作用在 $\sigma_0$ 生成的表示 $\pi$ 上。这就证明了 (a)。

回忆 $\sigma_0$ 是 $\tau$ 的约化中出现的权之一，于是我们把 $\tau$ 的约化中其余的权标为 $\sigma_1$、$\sigma_2$、$\sigma_3$。由于 $\tau$ 作为 $K$ 的表示有中心特征，每个 $\sigma_i$ 作为 $K$ 的表示也有中心特征，且这些特征都相同（它们就是 $\tau$ 的中心特征的约化）。特别地，$\pi$ 的中心特征与诸 $\sigma_i$ 的中心特征相容（即前者的限制到 $K\cap Z$ 与后者一致），因此我们可以并且确实把每个 $\sigma_i$ 延拓为 $KZ$ 的表示，让 $Z$ 通过 $\pi$ 的中心特征作用。

于是我们可以选取标号，使得对每个 $i=0,\dots,3$，都有一个嵌入（最初由 Pa\v{s}k\={u}nas 构造）
\begin{equation}\tag{7.7.53}
c\text{-}\mathrm{Ind}_{KZ}^G\sigma_{i+1}\oplus c\text{-}\mathrm{Ind}_{KZ}^G\sigma'_{i+1}\into c\text{-}\mathrm{Ind}_{KZ}^G\sigma_i/T,
\end{equation}

\ysrc{329}
其中 $i+1$ 按模 $4$ 计算，而 $\sigma'_{i+1}$ 是另一个 Serre 权，它不同构于任何 $\sigma_j$（$j=0,\dots,3$）。（关于这一点的稍详细说明，见下文注 7.7.54。）

以 $\pi'(\sigma_i)$ 记 (7.7.53) 的余核。一个基本事实是 $\pi'(\sigma_i)=\pi'(\sigma_i^{(s)})$，其中 $\sigma_i^{(s)}$ 记这样的 Serre 权：其最高权由 $\sigma_i$ 的最高权经 Weyl 群的对合 $s$ 作用而得。特别地，$\pi'(\sigma_i)$ 不仅是 $c\text{-}\mathrm{Ind}_{KZ}^G\sigma_i$ 的商，也是 $c\text{-}\mathrm{Ind}_{KZ}^G\sigma_i^{(s)}$ 的商。另一个基本事实是：诸 $\sigma_i^{(s)}$（$i=0,\dots,3$）中没有一个同构于任何 $\sigma_j$；也就是说，这些表示都不是 $\tau$ 的约化的组成因子。（同样见下文注 7.7.54。）由我们对 $\pi'$ 的 $K$-基座所作的假设，这意味着任何态射 $c\text{-}\mathrm{Ind}_{KZ}^G\sigma_i^{(s)}\ra\pi'$（对任何 $i$ 值）都必须为零，从而任何态射 $\pi'(\sigma_i)\ra\pi'$ 也必须为零。由于 $\pi\subseteq\pi'$，把 $\pi'$ 换作 $\pi$ 也有同样的结论。

现在考虑由包含 $\sigma_0\subseteq\pi$ 诱导的态射 $c\text{-}\mathrm{Ind}_{KZ}^G\sigma_0\ra\pi$。超奇异的假设蕴含（事实上，大致就意味着）该态射穿过 $c\text{-}\mathrm{Ind}_{KZ}^G\sigma_0/T$。由于它非零（其像含 $\sigma_0$），它必定不穿过 $\pi'(\sigma_i)$，因此它在 (7.7.53) 的定义域上不恒为零。然而它在该定义域的第二个直和项上必定为零（由我们对 $\pi'$ 的 $K$-基座所作的假设），因而它在第一个直和项上非零。于是我们得到一个非零态射 $c\text{-}\mathrm{Ind}_{KZ}^G\sigma_{i+1}\ra\pi$。等价地，我们得到一个嵌入 $\sigma_{i+1}\into\pi$。

如此继续，我们循环过所有权 $\sigma_i$，发现 $\pi'$ 在其 $K$-基座中含所有这些权，而且事实上它们都含于 $\pi$。这就完成了命题的证明。
\end{proof}

\begin{rem}{7.7.54}\label{rem:a4-7-7-54}
这里只想就我们上面用到的关于 Serre 权的事实略作补充：若以 $\chi_i$ 记 $\sigma_i$ 的最高权，则 $\mathrm{PS}(\chi_i)$ 的约化中的权为 $\sigma_i$、$\sigma_{i+1}$、$\sigma'_{i+1}$ 与 $\sigma_i^{(s)}$\tn{底本此处作 $\mathrm{PS}(\chi_i)$，无上横线；而 §7.7.50 中同一表示记作 $\ol{\mathrm{PS}}({\cdot})$。译文照录原文。}。而（充分一般的）尖点型的约化与（充分一般的）主级数型的约化或者没有公共权，或者恰有两个公共权。因此，既然 $\tau$ 的约化含 $\sigma_i$ 与 $\sigma_{i+1}$，它就不含 $\sigma'_{i+1}$ 或 $\sigma_i^{(s)}$。
\end{rem}

\begin{rem}{7.7.55}\label{rem:a4-7-7-55}
2006 年 AIM 与会者的原意，是取 $\pi'$ 为某条志村曲线（Shimura curve）的完备上同调中的 $\mathfrak m$-挠\tn{底本此处作 the completed cohomology of a Shimura curves，冠词与单复数均不顺；疑为 a Shimura curve 或 Shimura curves 之误。译文照录其意。}（其中 $\mathfrak m$ 记某个适当 Hecke 代数的极大理想），来应用命题 7.7.52。此时在适当的「一般性」假设下，命题的假设已由 [Gee11] 得知成立。更近期地，已经证明此时甚至有 $\pi'=\pi$（见 [BHH${}^{+}$21, Cor. 1.3.10]）。2006 年 AIM 会议后不久，Breuil 与 Pa\v{s}k\={u}nas [BP12] 直接构造出了命题 7.7.52 中的不可约容许表示 $\pi$。事实上，他们给出了无穷多个两两不同构的此类表示。
\end{rem}

命题 7.7.52 的证明表明，实际上有一个标量不变量附着于 $\pi'$——或者更确切地说，附着于它的不可约子表示 $\pi$。具体而言，若我们从包含 $\sigma_0\subseteq\pi$ 出发并循环四次，就回到一个嵌入 $\sigma_0\into\pi$；由我们的重数一假设，它必定等于原嵌入的一个非零标量倍。为使这一点良定义，必须把诸嵌入 (7.7.53) 固定下来，而稍加小心就可以做到

\ysrc{330}
（我们将在下文说明如何在应用 $\mathfrak A$ 之后做到这一点），于是得到 $\pi$ 的一个标量不变量。

更确切地说，取定非零标量 $\lambda=\lambda_0\in k^\times$，并令 $i=1,2,3$ 时 $\lambda_i=1$，考虑两步复形
\begin{equation}\tag{7.7.56}
\begin{tikzcd}[column sep=huge,row sep=small]
\displaystyle\bigoplus_{i=0}^{3}\bigl(c\text{-}\mathrm{Ind}_{KZ}^G\sigma_i\oplus c\text{-}\mathrm{Ind}_{KZ}^G\sigma'_i\bigr)\arrow[r] & \displaystyle\bigoplus_{i=0}^{3}c\text{-}\mathrm{Ind}_{KZ}^G\sigma_i/T,
\end{tikzcd}
\end{equation}
其中箭头在直和项 $c\text{-}\mathrm{Ind}_{KZ}^G\sigma_i\oplus c\text{-}\mathrm{Ind}_{KZ}^G\sigma'_i$ 上由
\[
\begin{aligned}
c\text{-}\mathrm{Ind}_{KZ}^G\sigma_i\oplus c\text{-}\mathrm{Ind}_{KZ}\sigma'_i\ni(v,w)\mapsto{}&(v+w,-\lambda_iv\bmod T)\\
&\in c\text{-}\mathrm{Ind}_{KZ}^G\sigma_{i-1}/T\oplus c\text{-}\mathrm{Ind}_{KZ}^G\sigma_i/T
\end{aligned}
\]
这里我们把 $v$ 与 $w$ 看作 $c\text{-}\mathrm{Ind}_{KZ}\sigma_{i-1}/T$ 的元素\tn{底本此处（及紧接其后的两处）作 $c\text{-}\mathrm{Ind}_{KZ}\sigma'_i$、$c\text{-}\mathrm{Ind}_{KZ}\sigma_{i-1}/T$，均缺上标 $G$；他处则作 $c\text{-}\mathrm{Ind}_{KZ}^G$。译文照录原文。}（经由 (7.7.53)）。

\begin{lem}{7.7.57}\label{lem:a4-7-7-57}
(7.7.56) 的箭头是单射。
\end{lem}

\begin{proof}
若核非零，则它有非零的 $K$-基座，于是对某个权 $\sigma'$，存在非零态射 $c\text{-}\mathrm{Ind}_{KZ}^G\sigma'\ra$ 该核。但若
\[
c\text{-}\mathrm{Ind}_{KZ}^G\sigma'\ra\bigoplus_{i=0}^{3}\bigl(c\text{-}\mathrm{Ind}_{KZ}^G\sigma_i\oplus c\text{-}\mathrm{Ind}_{KZ}^G\sigma'_i\bigr)
\]
非零，则 $\sigma'$ 等于某个 $\sigma_i$ 或 $\sigma'_i$（我们假设这些权充分一般），且该态射只是到相应直和项的嵌入。引理于是由下述事实得出：(7.7.56) 的箭头在其定义域的每个直和项上都是单射。
\end{proof}

以 $\Pi_\lambda$ 记 (7.7.56) 的余核。于是给出态射 $\Pi_\lambda\ra\pi$（$\pi$ 为某个其他表示）等价于给出 $\pi$ 中 $\sigma_0$ 的一个拷贝（或事实上任一 $\sigma_i$ 的拷贝），使得命题 7.7.52 的权循环在 $\pi$ 中「奏效」，且相应的标量等于 $\lambda$。

我们可以计算附着于 $\Pi_\lambda$ 的凝聚层（更确切地说，凝聚层复形）。答案依赖于 $\lambda$。事实上，几何使我们能把诸嵌入 (7.7.53) 标准化。具体而言，箭头 $c\text{-}\mathrm{Ind}_{KZ}^G\sigma_{i+1}\ra c\text{-}\mathrm{Ind}_{KZ}^G\sigma_i/T$ 与 $c\text{-}\mathrm{Ind}_{KZ}^G\sigma'_{i+1}\ra c\text{-}\mathrm{Ind}_{KZ}^G\sigma_i/T$ 在凝聚层一侧分别对应箭头 $\OO_{C_{\sigma_{i+1}}}(-1,0)\ra\OO_{C_\sigma}(-1)/(vy)$ 与 $\OO_{C_{\sigma'_{i+1}}}(-1,0)\ra\OO_{C_\sigma}(-1,0)/(vy)$\tn{底本此处前一箭头作 $\OO_{C_\sigma}(-1)/(vy)$，后一箭头作 $\OO_{C_\sigma}(-1,0)/(vy)$；按后者及上下文，前者疑脱 $0$。译文照录原文。}，而我们可以直接取它们为由包含 $C_\sigma^{vy=0}\into C_{\sigma_{i+1}},C_{\sigma'_{i+1}}$ 诱导的限制映射。对态射 (7.7.56) 应用 $\mathfrak A$，并按上述方式解释其中的诸态射，就可以（稍费些力气）显式算出复形 $\mathfrak A(\Pi_\lambda)$。

我们把这计算的结果总结为下述论断，它特别表明 $\lambda=1$ 这个值是特殊的。

\textbf{结论 7.7.58.}
\begin{enumerate}[(1)]
\item $H^0\bigl(\mathfrak A(\Pi_1)\bigr)$ 是支撑在 $x$ 处的摩天楼层（skyscraper），其中 $x$ 是各分支 $C_{\sigma_i}$（$i=0,\dots,3$）相交所得的闭点；而 $H^{-1}\bigl(\mathfrak A(\Pi_1)\bigr)$ 是支撑在 $\bigcup_{i=0}^{3}\bigl(C_{\sigma_i}\cup C_{\sigma'_i}\bigr)$ 上的线丛。
\item 若 $\lambda\neq1$，则 $H^0\bigl(\mathfrak A(\Pi_\lambda)\bigr)=0$，而 $H^{-1}\bigl(\mathfrak A(\Pi_\lambda)\bigr)$ 是支撑在 $C:=\bigcup_{i=0}^{3}\bigl(C_{\sigma_i}\cup C_{\sigma'_i}\bigr)$ 上的无挠层（torsion-free sheaf）。更确切地说，它是 $C\setminus\{x\}$ 上某个线丛推到 $C$ 上的像，该线丛依赖于参数 $\lambda$。
\end{enumerate}
\ysrc{331}
\subsubsection*{7.7.59.\quad 不可约商}
考虑短正合序列
\[
0\ra\mathcal K\ra\Pi_\lambda\ra\pi\ra0
\]
（其中 $\pi$ 是如上所说的不可约容许表示，而 $\mathcal K$ 定义为该满射的核）。过渡到（如今仅为假设性的）层，并注意 $\mathfrak A$ 的振幅为 $[-1,0]$（见注 6.1.28），我们得到长正合序列
\[
\begin{aligned}
0&\ra H^{-1}\bigl(\mathfrak A(\mathcal K)\bigr)\ra H^{-1}\bigl(\mathfrak A(\Pi_\lambda)\bigr)\ra H^{-1}\bigl(\mathfrak A(\pi)\bigr)\\
&\ra H^0\bigl(\mathfrak A(\mathcal K)\bigr)\ra H^0\bigl(\mathfrak A(\Pi_\lambda)\bigr)\ra H^0\bigl(\mathfrak A(\pi)\bigr)\ra0.
\end{aligned}
\]
若 $\lambda\neq1$，则我们得出 $H^0\bigl(\mathfrak A(\pi)\bigr)=0$（因为 $H^0\bigl(\mathfrak A(\Pi_\lambda)\bigr)=0$），从而 $\mathfrak A(\pi)$ 支撑在次数 $-1$ 上。若 $\lambda=1$，则由于 $H^0\bigl(\mathfrak A(\Pi_1)\bigr)$ 是支撑在 $x$ 处的摩天楼层，我们看到 $H^0\bigl(\mathfrak A(\pi)\bigr)$ 或者同样是支撑在 $x$ 处的摩天楼层，或者为 $0$。

另一方面，由于 $\pi$ 不是有限表现的 [Sch15b]，引理 A.4.6 (3) 蕴含 $H^{-1}\bigl(\mathfrak A(\pi)\bigr)$ 不凝聚，无论 $\lambda$ 取何值。

关于层 $H^{-1}\bigl(\mathfrak A(\pi)\bigr)$，我们能说些什么呢？它有一条性质如下：由于对 $\sigma'\neq\sigma_i$（$i=0,\dots,3$）有 $\Hom(c\text{-}\mathrm{Ind}_{KZ}^G\sigma',\pi)=0$\tn{底本此处作「if $\sigma'\neq\sigma_i$ for some $i=0,\dots,3$」，按文意当为「对一切 $i=0,\dots,3$ 皆有 $\sigma'\neq\sigma_i$」。译文照录原文。}，我们发现对任何这样的 $\sigma'$ 都有
\[
\Ext^1\bigl(\OO_{C_{\sigma'}},H^{-1}(\mathfrak A(\pi))\bigr)=0.
\]
若 $\lambda=1$，如上所说的表示 $\pi$（即 $\Pi_\lambda$ 的一个不可约商，其 $K$-基座恰由诸 $\sigma_i$ 构成、每个重数一）也可以通过整体方式构造，即经由定四元数代数与/或志村曲线的完备上同调。（在此情形下 $\lambda$ 必然等于 $1$，这已由 Dotto--Le [DL21] 证明。完备上同调中的表示已被许多作者研究过，其中 Taylor--Wiles 拼接是关键技术；关于已知结果的简要综述，我们请读者参看 [BHH${}^{+}$21] 的引言。）

当 $\pi$ 经由完备上同调构造时，函子 $\mathfrak A$ 与完备上同调之间所期待的关系提示，映射 $H^0\bigl(\mathfrak A(\Pi_1)\bigr)\ra H^0\bigl(\mathfrak A(\pi)\bigr)$ 应当非零，因此至少在此情形下，我们预期 $H^0\bigl(\mathfrak A(\pi)\bigr)$ 等于 $x$ 处的摩天楼层。我们在 §7.8 中进一步探讨层 $\mathfrak A(\pi)$：在那里我们将看到，我们的猜想蕴含不可约表示 $\pi$ 应当由条件 $H^0\bigl(\mathfrak A(\pi)\bigr)\neq0$ 唯一确定（见注 7.8.6）。

\subsection{伴随函子}
回到猜想 6.1.15 的情形，以及连续函子 $\mathfrak A:\Ind D^b_{\mathrm{f.p.}}(\mathrm{sm}\,G)\ra\Ind\Coh(\XX_d)$。由引理 A.4.6 (1)（即伴随函子定理），$\mathfrak A$ 有右伴随 $\mathfrak B:\Ind\Coh(\XX_d)\ra\Ind D^b_{\mathrm{f.p.}}(\mathrm{sm}\,G)$。由同一引理的第 (2) 部分（它适用，因为按假设 $\mathfrak A$ 保持紧对象；见注 6.1.29），我们看到 $\mathfrak B$ 也是连续的。

\begin{rem}{7.8.1}\label{rem:a4-7-8-1}
在注 1.4.6 的设定下（即 Fargues--Scholze 猜想的 $p$-进版本），函子 $\mathfrak A$ 与 $\mathfrak B$ 想必由函子 $i_!$ 与 $i^!$ 诱导，其中 $i$ 是开 stratum $[*/G(F)]\into\Bun_G$ 的包含。
\end{rem}

\ysrc{332}
我们可以用上一节的计算来研究 $\mathrm{PGL}_2(\Q_p)$ 与 $\mathrm{PGL}_2(\Q_{p^2})$ 情形下的函子 $\mathfrak B$。（当然，这并不完全是猜想 6.1.15 的设定，但上述讨论原封不动地适用于这一情形。）这将使我们能把本文的猜想与经由志村曲线上同调研究 $p$-进及 mod $p$ 局部 Langlands 对应的更传统途径联系起来，尤其与下述希望联系起来：可以把每个表示 $\overline{\rho}:\Gal_{\Q_{p^2}}\ra\mathrm{GL}_2(\F_p)$ 与 $\mathrm{GL}_2(\Q_{p^2})$ 的一个特定的容许表示联系起来，而该表示由它出现在与 $\overline{\rho}$ 的某个整体提升相关的任一（也就是说，每一个！）整体背景中而确定（例如见 [Bre10a, Q.\ 3.11]）。我们将在下文注 7.8.6 中回到这一主题。

我们可以借助 $\mathfrak A$ 与 Taylor--Wiles 拼接之间所期待的相容性来研究 $\mathfrak B$ 的预期性质。设 $F/\Q_p$ 任意，$\XX$ 为行列式为 $\varepsilon^{-1}$ 的秩 $2$ $(\phi,\Gamma)$-模构成的叠；并设 $L_\infty$ 为（当 $F\neq\Q_p$ 时是假设性的）pro-凝聚层，它是 $\mathfrak A$ 的核，即
\[
\mathfrak A(\pi):=L_\infty\otimes^L_{\OO[[G]]}\pi.
\]
则伴随函子由
\begin{equation}\tag{7.8.2}
\mathfrak B(\FF):=\mathrm{RHom}_{\mathrm{Pro-}\OO_\XX}(L_\infty,\FF)
\end{equation}
给出。

现在设 $x:\Spec k\ra\XX$ 是有限型点，对应连续表示 $\overline{\rho}:\Gal_F\ra\mathrm{GL}_2(k)$，并设 $f:\Spf R_\infty\ra\XX$ 是由在 $\overline{\rho}$ 的某个整体化处拼接而得的 versal 态射。由注 6.1.31，我们预期 $f^*L_\infty$ 等于拼接模 $M_\infty$。对一般的 $\overline{\rho}$，环 $R_\infty$ 形式光滑，而由 [GN22, Prop. 4.3.1] 以及 [BHH${}^{+}$23, HW22, Wan23] 的结果可得 $M_\infty$ 在 $R_\infty$ 上 pro-平坦（pro-flat）。我们预期它甚至是 pro-自由的。相应地，我们预期可以用下述公式计算 $\mathfrak B(\FF)$：
\[
\begin{aligned}
\mathfrak B(\FF)&=\mathrm{RHom}_{\mathrm{Pro-}\OO_\XX}(L_\infty,\FF)\\
&=R\Gamma\circ\mathcal Hom_{\mathrm{Pro-}\OO_\XX}(L_\infty,\FF)\\
&=R\Gamma\bigl(\XX,\mathcal Hom_{\mathrm{Pro-}\OO_\XX}(L_\infty,\FF)\bigr).
\end{aligned}
\]

记 $G_x=\Aut_{\Gal_F}(\overline{\rho})$。进一步假设 $\overline{\rho}^{\mathrm{ss}}$ 不是纯量（若如上假设 $\overline{\rho}$ 是一般的，则自动如此）；于是当 $\overline{\rho}$ 不可约或非分裂可约时 $G_x=\mu_2$，而当 $\overline{\rho}$ 是两个不同特征标之和时 $G_x=\mathbf G_m$。于是 $x$ 处的剩余带叠（residual gerbe）是仿射浸入\footnote{若 $\overline{\rho}$ 半单，则 $i$ 是闭浸入；否则 $i$ 分解为 $[\Spec k/\mu_2]\xrightarrow{\ \mathrm{open}\ }[\mathbf A^1_k/\mathbf G_m]\xrightarrow{\ \mathrm{closed}\ }\XX$，从而把它表示为仿射浸入的复合。}
\[
i:[\Spec k/G_x]\into\XX,
\]
并且可以构造拟凝聚层 $\delta_x:=i_*\OO_{[\Spec k/G_x]}$。若如上所述 $f:\Spf R_\infty\ra\XX$ 是由 Taylor--Wiles 拼接而得的 $x$ 处的 versal 态射，$\mathfrak m_\infty$ 记 $R_\infty$ 的极大理想，则我们看到
\[
\mathcal Hom_{\mathrm{Pro-}\OO_\XX}(L_\infty,\delta_x)=M_\infty^\vee[\mathfrak m_\infty]
\]
带有某个 $G_x$-作用。（要看出这一点，注意 $L_\infty\otimes_{\OO_\XX}\delta_x=M_\infty/\mathfrak m_\infty M_\infty$ 带有某个 $G_x$-作用，而 $M_\infty/\mathfrak m_\infty M_\infty$ 的连续 $k$-对偶等于 $M_\infty^\vee[\mathfrak m_\infty]$。）注意：按 Taylor--Wiles 拼接对 $M_\infty$ 的构造，$M_\infty^\vee[\mathfrak m_\infty]$ 等于完备上同调中的 $\mathfrak m$-挠，其中 $\mathfrak m$ 是某个

\ysrc{333}
适当 Hecke 代数的极大理想。由于 $G_x$ 线性约化（它等于 $\mu_2$ 或 $\mathbf G_m$ 之一），我们于是看到
\[
\mathfrak B(\delta_x)=\bigl(M_\infty^\vee[\mathfrak m_\infty]\bigr)^{G_x}.
\]
令 $n$ 变化，我们可以通过取值
\[
\mathfrak B\bigl(\delta_x(n)\bigr)=\bigl(M_\infty^\vee[\mathfrak m_\infty](n)\bigr)^{G_x}
\]
还原出整个 $M_\infty^\vee[\mathfrak m_\infty]$。或者，与其逐个考虑各种扭 $\delta_x(n)$，不如令 $g:\Spec k\ra\XX$ 记复合
\[
\Spec k\ra[\Spec k/G_x]\xrightarrow{\ i\ }\XX,
\]
于是我们可以对拟凝聚层 $\Delta_x:=g_*\OO_{\Spec k}$ 应用 $\mathfrak B$。（这是 $x$ 的剩余带叠上对应于 $G_x$ 的正则表示的那个拟凝聚层。）事实上，我们计算得
\[
\begin{aligned}
\mathfrak B(\Delta_x)&=\mathrm{RHom}_{\mathrm{Pro-}\OO_\XX}\bigl(L_\infty,g_*\OO_{\Spec\mathcal F}\bigr)\\
&=\mathrm{RHom}_{\mathrm{Pro-}\OO_{\Spec k}}\bigl(g^*L_\infty,\OO_{\Spec\mathcal F}\bigr)\\
&=\Hom_{\mathrm{cont}}\bigl(M_\infty/\mathfrak m_\infty M_\infty,k\bigr)=M_\infty^\vee[\mathfrak m_\infty].
\end{aligned}
\]
因此 $x\mapsto\mathfrak B(\Delta_x)$ 应当还原出上文提到的、所期待的「传统」mod $p$ 局部 Langlands 对应。由考虑层-Hom $\mathcal Hom_{\mathrm{Pro-}\OO_\XX}(L_\infty,\delta_x)$ 而来的 $M_\infty^\vee[\mathfrak m_\infty]$ 上的 $G_x$-作用，是 mod $p$ Langlands 对应中某种额外的结构；据我们所知，它此前未被考虑过（尽管可以说它隐含在 Breuil 及其合作者在 [Bre11] 与 [BH15] 等文中所表达的思想之中）。

\subsubsection*{7.8.3.\quad 一些例子}
在 $F=\Q_p$ 或 $F=\Q_{p^2}$ 的情形，我们对 $M_\infty^\vee[\mathfrak m]$ 有所了解（前一情形由 M.\,E.\ 的局部--整体相容性定理 [Eme11a]，后一情形由 [BHH${}^{+}$21]），并且对 $G_x$ 如何作用，或已知、或可猜测。我们有下列例子。

\begin{enumerate}[(1)]
\item 若 $F=\Q_p$ 且 $\overline{\rho}=\overline{\chi}_1\oplus\overline{\chi}_2$，并满足 $\overline{\chi}_1\overline{\chi}_2^{-1}\neq\mathrm{triv}$ 或 $\varepsilon^{\pm1}$，则 $M_\infty^\vee[\mathfrak m]$ 是两个不可约主级数表示 $\pi_1\oplus\pi_2$ 的直和，而且事实上，对 $\pi_1,\pi_2$ 的适当排序，有
\[
\mathcal Hom_{\mathrm{Pro-}\OO_\XX}(L_\infty,\delta_x)=\pi_1(1)\oplus\pi_2(-1)
\]
（其中扭表示 $\mathbf G_m$-作用的权；注意表观上的不对称性源于：为了把 $G_x$ 与 $\mathbf G_m$ 等同，需要给 $\overline{\chi}_1$、$\overline{\chi}_2$ 取一种排序）。在此情形 $\Delta_x:=\bigoplus_{n=-\infty}^{\infty}\delta_x(n)$，因此
\[
\mathfrak B(\Delta_x)=\mathfrak B\bigl(\delta_x(-1)\bigr)\oplus\mathfrak B\bigl(\delta_x(1)\bigr)=\pi_1\oplus\pi_2.
\]
若 $\sigma_1,\sigma_2$ 是 $\overline{\rho}$ 的两个 Serre 权（取定一种排序），则我们发现 $(\mathfrak A\circ\mathfrak B)\bigl(\delta_x(-1)\bigr)=\mathfrak A(\pi_1)$ 支撑在 $\XX(\sigma_1)$ 上，而 $(\mathfrak A\circ\mathfrak B)\bigl(\delta_x(1)\bigr)=\mathfrak A(\pi_2)$ 支撑在 $\XX(\sigma_2)$ 上。作为附带说明，我们指出：虽然 $\delta_x$ 的支撑在交 $\XX(\sigma_1)\cap\XX(\sigma_2)$ 上，但这些表示并非如此；表示 $\mathfrak A(\pi_i)$ 支撑在相应分支 $\XX(\sigma_i)$ 中对应于半单化为 $\overline{\rho}$ 的伽罗瓦表示的那个闭轨迹上。因此 $(\mathfrak A\circ\mathfrak B)(\Delta_x)=\mathfrak A(\pi_1)\oplus\mathfrak A(\pi_2)$ 支撑在这两条轨迹之并上，而后者正是所有以 $\overline{\rho}$ 为其半单化的伽罗瓦表示构成的轨迹。

\ysrc{334}
\item 若 $F=\Q_p$ 且 $\overline{\rho}$ 是特征标 $\chi_1$ 与 $\chi_2$ 的非分裂扩张，同样满足 $\chi_1\chi_2^{-1}\neq\mathrm{triv}$ 或 $\varepsilon^{\pm1}$，则 $M_\infty^\vee[\mathfrak m]$ 是例 (1) 中出现的那两个不可约主级数 $\pi_1$ 与 $\pi_2$ 的非分裂扩张，并且
\[
\mathcal Hom_{\mathrm{Pro-}\OO_\XX}(L_\infty,\delta_x)=M_\infty^\vee[\mathfrak m](1),
\]
其中的扭表明 $\mu_2$ 的作用非平凡。在此情形 $\Delta_x:=\delta_x\oplus\delta_x(1)$，因此
\[
\mathfrak B(\Delta_x)=\mathfrak B\bigl(\delta_x(1)\bigr)=M_\infty^\vee[\mathfrak m]
\]
（即 $\pi_1$ 与 $\pi_2$ 的那个非分裂扩张）。此外，$(\mathfrak A\circ\mathfrak B)(\Delta_x)=(\mathfrak A\circ\mathfrak B)\bigl(\delta_x(1)\bigr)$ 是 $\mathfrak A(\pi_1)$ 与 $\mathfrak A(\pi_2)$ 的非分裂扩张，特别地其支撑等于对应于半单化等于 $\chi_1\oplus\chi_2$ 的伽罗瓦表示的那个闭轨迹。

\item 若 $F=\Q_p$ 且 $\overline{\rho}$ 不可约，则 $M_\infty^\vee[\mathfrak m]=\pi$ 是一个超奇异不可约表示，并且
\[
\mathcal Hom_{\mathrm{Pro-}\OO_\XX}(L_\infty,\delta_x)=\pi(1)
\]
（其中的扭表明 $\mu_2$ 的作用非平凡）。与例 (2) 一样，我们有 $\Delta_x:=\delta_x\oplus\delta_x(1)$，因此
\[
\mathfrak B(\Delta_x)=\mathfrak B\bigl(\delta_x(1)\bigr)=\pi,
\]
并且 $(\mathfrak A\circ\mathfrak B)(\Delta_x)=(\mathfrak A\circ\mathfrak B)\bigl(\delta_x(1)\bigr)=\delta_x(1)$。

\item 若 $F=\Q_{p^2}$ 且 $\overline{\rho}=\chi_1\oplus\chi_2$，其中 $\chi_1,\chi_2$ 适当一般，则由 [BHH${}^{+}$21, Cor. 1.3.10] 有 $M_\infty^\vee[\mathfrak m]=\pi_1\oplus\pi\oplus\pi_2$，其中 $\pi_1,\pi_2$ 是显式的不可约主级数表示，而 $\pi$ 是（非显式的！）超奇异不可约表示。我们预期
\[
\mathcal Hom_{\mathrm{Pro-}\OO_\XX}(L_\infty,\delta_x)=\pi_1(2)\oplus\pi\oplus\pi_2(-2).
\]
假设此式成立，我们看到
\[
\mathfrak B(\Delta_x)=\mathfrak B\bigl(\delta_x(-2)\bigr)\oplus\mathfrak B(\delta_x)\oplus\mathfrak B\bigl(\delta_x(2)\bigr)=\pi_1\oplus\pi\oplus\pi_2.
\]
（这里我们用到（如例 (1) 中那样）$\Delta_x:=\bigoplus_{n=-\infty}^{\infty}\delta_x(n)$。）特别地，我们发现 $\mathfrak B(\delta_x)=\pi$ 不可约且超奇异。（对特征标的非分裂扩张情形，可以用 [HW22] 的结果作类似的分析。）

\item 若 $F=\Q_{p^2}$ 且 $\overline{\rho}$ 不可约且适当一般，则由 [BHH${}^{+}$21, Thm. 1.3.11] 有 $M_\infty^\vee[\mathfrak m]=\pi$，其中 $\pi$ 是超奇异不可约表示。我们预期
\[
\mathcal Hom_{\mathrm{Pro-}\OO_\XX}(L_\infty,\delta_x)=\pi
\]
（$\mu_2$ 的作用平凡）。假设此式成立，我们看到
\[
\mathfrak B(\Delta_x)=\mathfrak B(\delta_x)=\pi
\]
不可约且超奇异。（这里我们用到（如例 (3) 中那样）$\Delta_x:=\delta_x\oplus\delta_x(1)$。）
\end{enumerate}

\begin{rem}{7.8.4}\label{rem:a4-7-8-4}
上文例 (4) 与 (5) 中的扭最初是由启发式得到的，依据是例 (1)、(2)、(3) 以及 Breuil 的论文 [Bre11] 所提示的哲学。然而，人们可以确认它们（在我们的猜想框架内，即假设 $\mathfrak A$ 存在且具有所猜想的性质，并且对充分一般的 Serre 权 $\sigma$，$\mathfrak A\bigl(c\text{-}\mathrm{Ind}_{KZ}^G\sigma\bigr)$ 的取值
% 注 7.8.4 的正文在 ch-a4-15.tex 开头（印刷页 335 首段，"……（如期望 7.6.12 所预言）的方式如下："）续完。
% 下方的 \end{rem} 为保持本片环境配平而暂加；合稿时应把它移到该续段之后（即 ch-a4-15 的 \begin{rem}{7.8.5} 之前）。
% 注 7.8.4 在此跨片断开，\end{rem} 见 ch-a4-15.tex 中 \begin{rem}{7.8.5} 之前。


% 本片止于印刷页 334（底本 PDF 347 页）末；上句在底本 PDF 348 页（印刷页 335）继续：
% "……（如期望 7.6.12 所预言）的方式如下：…"。续片 ch-a4-15.tex 自 PDF 348 页起接。
